\section{Sampling proposal implementation details}
\label{app:sampling-details}

\paragraph{Gradient features.}
We use gradient-informed proposals on the discrete graph space with target density
\(\pi_\theta(x)\propto \exp\left(-V_\theta(x)\right)\). We allow an effective inverse temperature
\(\beta_{\mathrm{mh}}\) in the \gls{mh} ratio below (\(\beta_{\mathrm{mh}}=1\) for the base model). We denote
the gradients with respect to real-valued one-hot node and edge features as follows.
\begin{align}
g^\mathcal{V} & := \operatorname{reshape}_{n,l_{\mathrm{node}}}\big(\nabla_{\hat{x}_\mathcal{V}} V_\theta(\hat{x})\big) \in \mathbb{R}^{n\times l_{\mathrm{node}}}\nonumber\\
g^\mathcal{E} & := \operatorname{reshape}_{\binom{n}{2},l_{\mathrm{edge}}}\big(\nabla_{\hat{x}_\mathcal{E}} V_\theta(\hat{x})\big) \in \mathbb{R}^{\binom{n}{2}\times l_{\mathrm{edge}}}\nonumber
\end{align}
Since we are working with undirected graphs, we use symmetrized edge gradients
\(g_{ij}=g_{ji}=\tfrac{1}{2}(g_{ij}+g_{ji})\) from the graph transformer.

\paragraph{Factorized proposal (multi-site).}
We sample all nodes and undirected edges independently from categorical logits
\begin{align}
\ell^\mathcal{V}_{i,c_i} & = \beta_{\mathrm{prop}}\big(g^\mathcal{V}_{i,x_i}-g^\mathcal{V}_{i,c_i}\big)-\lambda_\mathcal{V}\,\mathbb{I}_{\{c_i \neq x_i\}};~\forall~c_i \in \mathcal{X}_\mathrm{node}, \nonumber \\
\ell^\mathcal{E}_{{ij},c_{ij}} &= \beta_{\mathrm{prop}}\big(g^\mathcal{E}_{ij,x_{ij}}-g^\mathcal{E}_{ij,c_{ij}}\big)
-\lambda_\mathcal{E}\,\mathbb{I}_{\{c_{ij}\ne x_{ij}\}};~\forall~c_{ij} \in \mathcal{X}_\mathrm{edge} \nonumber
\end{align}
so the full proposal factorizes as
\begin{equation}
\begin{aligned}
q_{\beta_{\mathrm{prop}}}(x\to y)
&=
\prod_{i\in\mathcal{V}}\mathrm{Cat}\left(y_i;\mathrm{softmax}(\ell_i^\mathcal{V})\right)\\
&\quad\times\prod_{i<j}\mathrm{Cat}\left(y_{ij};\mathrm{softmax}(\ell_{ij}^\mathcal{E})\right).
\end{aligned}
\end{equation}
The current class has logit \(0\) (the \textit{stay} option), while \(\lambda_\mathcal{V},\lambda_\mathcal{E}\)
penalize changes.

\paragraph{Metropolis--Hastings refinement (single beta).}
For the standard Langevin proposal we set \(\beta_{\mathrm{prop}}=\beta_{\mathrm{mh}}=\beta\) and accept
\begin{equation}
\begin{aligned}
\alpha(x,y)
&=\min\{1,\exp(\Delta(x,y))\},\\
\Delta(x,y)
&=-\beta_{\mathrm{mh}}\big(V_\theta(y)-V_\theta(x)\big)
+\log\frac{q_{\beta_{\mathrm{prop}}}(y\to x)}{q_{\beta_{\mathrm{prop}}}(x\to y)}.
\end{aligned}
\end{equation}

\paragraph{Two-betas simulated annealing.}
We decouple the proposal temperature and target temperature by using
\(\beta_{\mathrm{prop}}\) in the logits above, while \(\beta_{\mathrm{mh}}\) controls the \gls{mh} ratio.
Optionally, \(\beta_{\mathrm{mh}}\) is annealed across \gls{mcmc} steps:
\begin{equation}
\beta_{\mathrm{mh}}(s)=
\begin{cases}
\beta_{\mathrm{mh}}^{\mathrm{init}}+\Delta\beta\,\dfrac{s}{S_{\mathrm{anneal}}-1}, & 0\le s<S_{\mathrm{anneal}},\\[6pt]
\beta_{\mathrm{mh}}^{\mathrm{final}}, & s\ge S_{\mathrm{anneal}}.
\end{cases}
\end{equation}
where \(\Delta\beta=\beta_{\mathrm{mh}}^{\mathrm{final}}-\beta_{\mathrm{mh}}^{\mathrm{init}}\).

\paragraph{Do-nothing resampling and proposal diffusion.}
If the independently sampled proposal equals the current state (all sites \textit{stay}),
we resample after softening the logits to increase randomness:
\(\beta_{\mathrm{prop}}\leftarrow\rho\,\beta_{\mathrm{prop}}\),
\(\lambda_\mathcal{V}\leftarrow\rho\,\lambda_\mathcal{V}\),
\(\lambda_\mathcal{E}\leftarrow\rho\,\lambda_\mathcal{E}\) with \(0<\rho<1\).
We repeat this a bounded number of times; if no change occurs, we keep the current
state and accumulate the \textit{stay} probability in \(\log q\). The reverse
transition uses the same resampling rule, preserving the \gls{mh} ratio.
